\newpage
\part*{ВВЕДЕНИЕ}
    Параметрическая идентификация -- определение неизвестных параметров математической модели объекта по результатам наблюдения его входа и выхода. Структура модели (передаточная функция) считается известной, а часть её коэффициентов -- нет.

    В работе идентификация моделируется на вычислительном эксперименте: истинные значения параметров известны заранее. Это позволяет оценить, насколько точно их удаётся восстановить. Работа состоит из следующих этапов:
    \begin{enumerate}
        \item по заданной передаточной функции составляется система дифференциальных уравнений и численно находится реакция объекта на ступенчатое воздействие $y^T(t)$; решение проверяется по точной формуле;
        \item реализуется метод оптимизации Гаусса-Зейделя ($GZ_1$) и проверяется на тестовых функциях с известным минимумом;
        \item к $y^T$ добавляется случайный шум трёх уровней, так получаются <<экспериментальные>> данные $y^{\text{Э}}$; проверяется генератор случайных чисел;
        \item для каждого опыта минимизируется целевая функция -- сумма квадратов отклонений модели от эксперимента. Найденные значения неизвестных параметров сравниваются с истинными, и анализируется влияние уровня шума на точность.
    \end{enumerate}

    Все вычисления выполняются программой на языке Fortran (см. приложение); все числа и графики в отчёте построены по её результатам.
